\section{Background and related literature}
\label{sec:background}

\subsection{Measuring the informal economy}

The size of the informal or shadow economy cannot be observed directly, and the literature has developed several indirect approaches \citep{SchneiderEnste2000, MedinaSchneider2018}. Currency-demand methods attribute part of the demand for cash to transactions that avoid taxation or registration. For Colombia, \citet{ArangoMisasLopez2006} estimate the underground economy from 1976 to 2003 with a cash-demand function in state-space form. Multiple-indicators multiple-causes (MIMIC) models treat the shadow economy as a latent variable linked to causes such as taxation and regulation and to indicators such as currency and labor participation \citep{MedinaSchneider2018}. Panels for Colombia and most other countries are available from \citet{SchneiderBuehnMontenegro2010}, \citet{MedinaSchneider2019} and, with an enhanced MIMIC, \citet{AsllaniDellAnnoSchneider2024}; \citet{Vuletin2008} estimates the informal economy of Latin America and the Caribbean with a structural equation model, and the World Bank database of \citet{ElginKoseOhnsorgeYu2021} collects DGE and MIMIC estimates. These methods produce a time path but need an external value to fix its level, and their structural interpretation is limited.

Model-based approaches use the equilibrium conditions of an economy with a formal and an informal sector. \citet{IhrigMoe2004} and \citet{BusatoChiarini2004} study the allocation of labor between a taxed formal sector and an untaxed informal one. \citet{ElginOztunali2012} calibrate a two-sector dynamic general equilibrium model to produce shadow-economy estimates for many countries, fixing the level in a base year with an external estimate. \citet{SolisGarciaXie2017} show that the growth rate imposed on shadow productivity in that approach is not consistent with the model's balanced-growth conditions and derive it from the observed trends instead; the level is again fixed with an external base-year value. \citet{OrsiRaggiTurino2014} estimate an underground economy for Italy in a Bayesian DSGE model with limited tax enforcement; observed consumption, investment, tax revenues and regular earnings identify the latent sector, and the prior is centered on the official estimate of the statistical office. They note that a key parameter of the underground technology has a posterior almost identical to its prior.

\subsection{Informality in business-cycle models}

A second literature uses informality to explain business-cycle facts in emerging economies. \citet{RestrepoEchavarria2014} shows that if national accounts measure informal output incompletely, measured consumption is more volatile than measured output. \citet{FernandezMeza2015} estimate a model of informal employment for Mexico and find that informality is countercyclical and amplifies the response to productivity shocks. The stochastic trend of the present model follows \citet{AguiarGopinath2007}, who show that permanent productivity shocks account for much of the cycle in emerging economies.

\subsection{Identification}

Identification in DSGE models has been studied in general by \citet{CanovaSala2009}, \citet{Iskrev2010} and \citet{KoopPesaranSmith2013}. Bayesian estimation always delivers a posterior, so a parameter that the likelihood does not inform is reported with a posterior close to its prior. The present paper applies three standard checks to the parameter that matters for the measurement of the informal economy: the ratio of posterior to prior dispersion, the profile likelihood, and the recovery of a known value in simulated data.

\subsection{Official statistics in Colombia}

DANE publishes quarterly national accounts in volumes chained to the 2015 base. Its national accounts include an imputation for the non-observed economy, estimated at 22.4\% of total production in the 2015 base \citep{DANE_SIECI}; this concept refers to gross production, not value added, and includes underground, illegal and informal activities and statistical corrections. Labor informality is measured by the household survey (GEIH) and has changed definition over time; the ICML project reconstructs consistent national and territorial series from 1993 \citep{VeraJaramillo2025ICML}.
